%! TeX program = lualatex
\documentclass[12pt]{article}

\usepackage{cmap}
\usepackage{newfloat}
\usepackage{tabularx}
\usepackage[newfloat]{minted}
\usepackage[font=small]{caption}
% \usepackage{indentfirst}
\usepackage{booktabs} % \toprule, \midrule, \bottomrule
\usepackage[table]{xcolor}

\usepackage{tikz}
\usetikzlibrary{shapes, arrows.meta, positioning, shadows}

\setminted{
    linenos,                % показывать номера строк
    breaklines,             % переносить длинные строки
    breakanywhere=true,     % переносить в любом месте
    tabsize=2,              % размер табуляции
    fontsize=\footnotesize,        % размер шрифта
    frame=lines,            % рамка сверху и снизу
    framesep=2mm,           % расстояние текста от рамки
    baselinestretch=1.1,    % межстрочный интервал
    autogobble=true,        % автоудаление общего отступа
    % xleftmargin=10pt,       % отступ слева
    % numbersep=5pt,          % расстояние между номерами и кодом
}

\usepackage[english, russian]{babel}

\usepackage[nopatch=footnote]{microtype}

\usepackage{float}
\usepackage{fontspec}
\usepackage[pdfborder={0 0 0}]{hyperref}

\usepackage{enumitem}

\makeatletter
\renewcommand\@seccntformat[1]{%
  \csname the#1\endcsname.\quad
}
\makeatother

\usepackage{tocloft}
\renewcommand{\cftsecaftersnum}{.}
\renewcommand{\cftsubsecaftersnum}{.}
\renewcommand{\cftsubsubsecaftersnum}{.}

\setmainfont{TimesNewerRoman}[
  Extension = .otf,
  Path = /nix/store/ihdqrn4p54awd89vrday9k53s8i47bbv-times-newer-roman-unstable-2018-09-11/share/fonts/opentype/,
  UprightFont = *-Regular,
  BoldFont = *-Bold,
  ItalicFont = *-Italic
]

\setmonofont{Ubuntu Mono}

\newcommand{\icon}[1]{\fontspec{UbuntuNerdFont}[Extension = .ttf,
  Path = /nix/store/h721jafy2n74x4k5p0hxbz722h0rncmx-nerd-fonts-ubuntu-3.3.0+0.83/share/fonts/truetype/NerdFonts/Ubuntu/,
  UprightFont = *-Regular,
BoldFont = *-Bold] #1}

\newcommand{\iicon}[1]{{\icon{#1}}}

\usepackage{graphicx}
\usepackage{fontawesome5} % Для иконок

\usepackage[left=2.0cm, right=2.0cm, top=1.0cm, bottom=1.0cm, includeheadfoot]{geometry}

% \usepackage[mocha, textcolor=true, pagecolor=true]{catppuccinpalette} \usemintedstyle{catppuccin-mocha}
\usepackage[latte, textcolor=true, pagecolor=true]{catppuccinpalette} \usemintedstyle{catppuccin-latte}

\usepackage{setspace}
\onehalfspacing

\usepackage{fancyhdr}
\fancyhf{}

\renewcommand{\sectionmark}[1]{\markboth{#1}{}}
\renewcommand{\subsectionmark}[1]{\markright{#1}}

% Настройка колонтитулов
\fancyhead[RO]{\rightmark}
\fancyhead[LO]{\leftmark}

\fancyfoot[L]{\hspace{8pt} \rule{\textwidth}{0.4pt}}
\fancyfoot[R]{\LARGE{\thepage}}
\fancyfoot[C]{}

\newcommand{\lablogo}
{
\begin{center}
    \huge{\textbf{Лабораторная работа №5}} \\
\end{center}
}

\newcommand{\colorURL}[1]{\textcolor{CtpBlue}{#1}}
\newcommand{\colorGIT}[1]{\textcolor{CtpLavender}{#1}}
\renewcommand{\texttt}[1]{{\small\ttfamily #1}}

\setlength{\headheight}{15.2pt}

% \DeclareCaptionLabelFormat{gostfigure}{Рисунок #2}
% \captionsetup[figure]{labelformat=gostfigure, labelsep=endash}

\DeclareCaptionLabelFormat{gostfigure}{Таблица #2}
\captionsetup[table]{labelformat=gostfigure, labelsep=endash}

\DeclareCaptionLabelFormat{gostlisting}{Листинг #2}
\newenvironment{code}
  {\par\captionsetup{type=listing}\noindent}
  {\par\vspace{\baselineskip}}
\captionsetup[listing]{labelformat=gostlisting, labelsep=endash}
% \DeclareFloatingEnvironment[
%   listname=Листинг,
% ]{listing}

\usepackage{amsmath}
\numberwithin{listing}{section}
\numberwithin{figure}{section}

\usepackage{tocloft}
\renewcommand{\cftsecleader}{\cftdotfill{\cftdotsep}}

\begin{document}
\pagestyle{empty}
\lablogo
\begin{center}
	\Large{Интеграция WPF-приложения с ASP.NET Core Web API.}
\end{center}
\vspace{1em}
\tableofcontents
\newpage
\listoflistings
\newpage

\pagestyle{fancy}

\section{Цель работы~\texorpdfstring{\iicon{}}{}}
\textbf{Цель:} организовать взаимодействие между клиентским WPF-приложением и серверным REST API, обеспечив обмен данными через HTTP-запросы.

% \begin{table}[H]
% 	\centering
% 	\rowcolors{2}{CtpBase}{CtpSurface0}
% 	\resizebox{\linewidth}{!}{
% 		% \label{tab:project_structure}
% 		\begin{tabular}{l|l}
% 			\toprule
% 			\textbf{Состав проекта}  & \textbf{Назначение}                                        \\
% 			\midrule
% 			StoreManagerApi          & ASP.NET Web API (серверная часть, отвечает за работу с БД) \\
% 			StoreManager\_4lab (WPF) & WPF-приложение (клиентская часть, UI для пользователя)     \\
% 			SQLite БД                & База данных, расположенная в store.db                      \\
% 			\bottomrule
% 		\end{tabular}
% 	}
% 	\caption{Структура проекта}
% \end{table}

\vspace{12pt}

\begin{figure}[h]
	\centering
	\begin{tikzpicture}[
		node distance=1cm,
		every node/.style={
				font=\small,
				align=center
			},
		block/.style={
				rectangle,
				rounded corners=3pt,
				draw=CtpBlue,
				fill=CtpMantle,
				drop shadow={opacity=0.2},
				minimum width=3.5cm,
				minimum height=1.2cm,
				text width=3.2cm,
				inner sep=5pt
			},
		arrow/.style={
		-{Stealth[length=3mm]},
		thick,
		CtpBlue
		},
		% title/.style={
		% 		rectangle,
		% 		fill=CtpBlue,
		% 		text=CtpBase,
		% 		font=\bfseries,
		% 		minimum width=8cm,
		% 		inner sep=8pt
		% 	}
		]

		% \node[title] (title) {Структура проекта};

		\node[block] (api) {StoreManagerApi \\ \footnotesize ASP.NET Web API};
		\node[block, left=3.2cm of api] (wpf) {StoreManager\_4lab \\ \footnotesize WPF-приложение};
		\node[block, right=3.2cm of api] (db) {SQLite БД \\ \footnotesize store.db};

		\draw[arrow] (wpf) -- node[below, sloped] {HTTP запросы} (api);
		\draw[arrow] (api) -- node[below, sloped] {Данные} (db);
		\draw[arrow] (db) -- node[above, sloped] {Результаты} (api);
		\draw[arrow] (api) -- node[above, sloped] {Ответы} (wpf);

		\node[below=1cm of api, text width=8cm, align=center] (desc) {
			\begin{tabular}{l|c}
				{\textbf{Компонент}} & {\textbf{Назначение}}        \\
				\midrule
				StoreManagerApi      & Серверная часть, работа с БД \\
				StoreManager\_4lab   & Клиентское WPF-приложение    \\
				SQLite БД            & Хранение данных в store.db
			\end{tabular}
		};

	\end{tikzpicture}
	\caption{Структура проекта}
\end{figure}

\noindent \textcolor{CtpRed}{Внимание:} проект ASP.net созданный в лабораторной №5 изменять в ходе данной лабораторной не нужно будет. Необходимо привести проект созданный в лабораторной №4 к тому виду, который будет рассмотрен в данной лабораторной и убедиться, что все функции (кнопки) выполняются без ошибок.

\newpage

\section{Общая информация~\texorpdfstring{\iicon{}}{}}
\begin{enumerate}
	\item WPF-приложение — это клиент, т.е. пользовательский интерфейс, установленный на компьютере пользователя.
	\item ASP.NET Web API — это сервер, который обрабатывает запросы, работает с базой данных и отправляет данные клиенту (например, в формате JSON).
\end{enumerate}

\subsection{Роли каждого проекта}
\noindent \textbf{ASP.NET (сервер):}
\begin{enumerate}
	\item Отвечает за доступ к базе данных.
	\item Использует DbContext из Entity Framework, чтобы читать и сохранять данные.
	\item Предоставляет HTTP API (например, GET, POST, PUT, DELETE), к которым можно обращаться через интернет или локальную сеть.
\end{enumerate}

\noindent \textbf{WPF (клиент):}
\begin{enumerate}
	\item Отображает данные пользователю.
	\item Отправляет HTTP-запросы к ASP.NET API, чтобы получить или отправить данные.
	\item Не подключается к базе данных напрямую, а работает через API. Поэтому DbContext здесь не нужен.
\end{enumerate}


\begin{table}[H]
	\centering
	\small
	\rowcolors{2}{CtpBase}{CtpSurface0}
	\begin{tabularx}{\linewidth}{p{4.2cm}|>{\centering\arraybackslash}X|>{\centering\arraybackslash}X}
		\toprule
		\textbf{Компонент}                     & \textbf{Лабораторная №4 (локальная БД)}                     & \textbf{Лабораторная №6 (через API)}              \\
		\midrule
		Хранение данных                        & SQLite через Entity\-Framework\-Core прямо в WPF-приложении & ASP.NET Core Web API, клиент обращается к серверу \\
		Репозитории                            & Прямой доступ к DbContext                                   & Используются Http\-Client и Api\-Service          \\
		Работа с данными                       & StoreDbContext и LI\-NQ-запросы                             & REST-запросы: GET, POST, PUT, DELETE              \\
		Добавление / удаление / редактирование & Прямое обращение к базе через EF                            & Отправка данных в контроллеры Web API             \\
		Обновление интерфейса                  & Автоматически при изменении коллекций                       & После получения данных через HTTP-запрос          \\
		\bottomrule
	\end{tabularx}
	\caption{Архитектурные изменения}
\end{table}


\subsection{Изменения в коде}
\textbf{Контекст базы данных (StoreDbContext)}
\begin{enumerate}
	\item Было: EF Core DbContext создавался и использовался прямо в MainViewModel и других ViewModel.
	\item Стало: База данных вынесена на сервер в ASP.NET Core Web API, клиент WPF не работает с EF напрямую.
\end{enumerate}
Причина изменения: разделение ответственности — клиент занимается только отображением, сервер — бизнес-логикой и хранением.

\noindent В данной таблице приведены ключевые различия в подходе к работе с данными между локальным доступом через DbContext (в лабораторной №4) и использованием Web API через HttpClient (в лабораторной №6).


\renewcommand{\texttt}[1]{{\footnotesize\ttfamily #1}}
\begin{table}[H]
	\centering
	\footnotesize
	\rowcolors{2}{CtpBase}{CtpSurface0}
	\begin{tabularx}{\linewidth}{X|>{\centering\arraybackslash}X|>{\centering\arraybackslash}X|>{\centering\arraybackslash}X}
		\toprule
		\textbf{Описание операции} & \textbf{Лабораторная №4 (локально)}             & \textbf{Лабораторная №6 (через Web API)}                    & \textbf{Изменение}              \\
		\midrule
		Добавление товара          & \texttt{\_con\-text.\-Products.Add\-(product);} & \texttt{await\_api.Add\-Product\-Async\-(product);}         & Используется API вместо EF Core \\
		Сохранение изменений       & \texttt{\_context.SaveChanges();}               & Выполняется внутри метода API                               & Уходит внутрь ASP.NET API       \\
		Загрузка товаров           & \texttt{foreach(varpin\_con\-text.\-Products)}  & \texttt{varproducts=await\_\-api.\-Get\-Products\-Async();} & Получение данных по HTTP        \\
		Создание заказа            & \texttt{\_context.Orders.Add\-(order);}         & \texttt{await\_api.AddOrder\-Async\-(order);}               & Заказ создается через API       \\
		Удаление заказа/товара     & \texttt{\_context.Orders.Remove\-(order);}      & \texttt{await \_api.DeleteOrder\-Async\-(order.Id);}        & Операция выполняется на сервере \\
		\bottomrule
	\end{tabularx}
	\caption{Основные изменения в методах используемых в приложении}
\end{table}

\renewcommand{\texttt}[1]{{\small\ttfamily #1}}

\noindent Таким образом, ключевым отличием лабораторной №6 стало отделение клиентской логики (WPF) от серверной (ASP.NET Core Web API) и переход от прямого доступа к базе данных к REST-взаимодействию через HTTP-запросы.

\newpage

\section{Реализация клиентского сервиса (\texttt{ApiService.cs})~\texorpdfstring{\iicon{}}{}}
\texttt{ApiService.cs} — клиентский сервис доступа к Web API Класс \texttt{ApiService} реализует клиентскую часть взаимодействия WPF-приложения с серверным Web API. Он инкапсулирует вызовы HTTP-запросов к контроллерам товаров и заказов, используя \texttt{HttpClient}. Это позволяет централизованно управлять всеми запросами к серверу и облегчает сопровождение кода.

\begin{code}
	\begin{minted}{csharp}
namespace StoreManager_6lab.Services
{
    public class ApiService
    {
        private readonly HttpClient _http; // HTTP-клиент для запросов к API

        public ApiService()
        {
            var handler = new HttpClientHandler
            {
                ServerCertificateCustomValidationCallback = (message, cert, chain, errors) => true // Разрешить самоподписанные сертификаты
            };

            _http = new HttpClient(handler)
            {
                BaseAddress = TryUseHttps()
                    ? new Uri("https://localhost:7259/api/")
                    : new Uri("http://localhost:5107/api/") // Базовый адрес сервера API
            };
        }

        // Попытка проверить доступность HTTPS API
        private bool TryUseHttps()
        {
            try
            {
                using var testClient = new HttpClient(
                    new HttpClientHandler { ServerCertificateCustomValidationCallback = (a, b, c, d) => true });
                var response = testClient.GetAsync("https://localhost:7259/api/Product").Result;
                return response.IsSuccessStatusCode;
            }
            catch
            {
                return false;
            }
        }

        // === Методы работы с товарами ===

        public async Task<List<Product>> GetProductsAsync() // Получить список товаров
        {
            try
            {
                return await _http.GetFromJsonAsync<List<Product>>("Product") ?? new List<Product>();
            }
            catch (Exception ex)
            {
                MessageBox.Show($"Ошибка при получении данных: {ex.Message}");
                return new List<Product>();
            }
        }

        public async Task AddProductAsync(Product product) => // Добавить новый товар
            await _http.PostAsJsonAsync("Product", product);

        public async Task UpdateProductAsync(Product product) => // Обновить товар по ID
            await _http.PutAsJsonAsync($"Product/{product.Id}", product);

        public async Task DeleteProductAsync(int id) => // Удалить товар по ID
            await _http.DeleteAsync($"Product/{id}");

        // === Методы работы с заказами ===

        public async Task<List<Order>> GetOrdersAsync() // Получить список заказов
        {
            try
            {
                return await _http.GetFromJsonAsync<List<Order>>("Order") ?? new List<Order>();
            }
            catch (Exception ex)
            {
                MessageBox.Show($"Ошибка при получении заказов: {ex.Message}");
                return new List<Order>();
            }
        }

        public async Task AddOrderAsync(Order order) => // Добавить новый заказ
            await _http.PostAsJsonAsync("Order", order);

        public async Task DeleteOrderAsync(int id) => // Удалить заказ по ID
            await _http.DeleteAsync($"Order/{id}");
    }
}
\end{minted}
	\captionof{listing}{Клиентского \colorGIT{\href{https://github.com/WebMasterIT/Csharp_Labs/blob/main/6lab/StoreManager_6lab/Services/AppService.cs}{сервиса}}}
	\label{lst:api_service}
\end{code}

\noindent Описание методов (к листингу~\ref{lst:api_service}):
\begin{enumerate}
	\item \texttt{GetProductsAsync()} — выполняет GET-запрос на получение всех товаров.
	\item \texttt{AddProductAsync(product)} — отправляет POST-запрос на добавление нового товара.
	\item \texttt{UpdateProductAsync(product)} — обновляет данные существующего товара по ID.
	\item \texttt{DeleteProductAsync(id)} — удаляет товар по указанному ID.
	\item \texttt{GetOrdersAsync()} — получает все заказы, включая их состав.
	\item \texttt{AddOrderAsync(order)} — отправляет заказ на сервер для добавления в БД.
	\item \texttt{DeleteOrderAsync(id)} — удаляет заказ по ID.
\end{enumerate}

\newpage

\section{Реализация \texttt{MainViewModel.cs}~\texorpdfstring{\iicon{}}{}}
\texttt{MainViewModel.cs} — основная ViewModel клиентского WPF-приложения Класс \texttt{MainViewModel} реализует основную бизнес-логику WPF-клиента. Он управляет товарами, корзиной, заказами и взаимодействует с \texttt{ApiService} для выполнения операций через Web API. Все команды пользовательского интерфейса (добавление, удаление, обновление) реализованы через команды MVVM-паттерна. В листинге~\ref{lst:main_view_model} представлены методы основной ViewModel разрабатываемого сервиса.
\begin{code}
	\begin{minted}{csharp}
public class MainViewModel : INotifyPropertyChanged
{
    private readonly ApiService _api = new(); // Экземпляр сервиса для обращения к Web API

    public ObservableCollection<Product> Products { get; set; } = new(); // Список всех товаров
    public ObservableCollection<Order> Orders { get; set; } = new(); // Список заказов
    public ObservableCollection<CartItem> CartItems { get; set; } = new(); // Список товаров в корзине

    // Вводимые пользователем данные о товаре
    public string NewProductName { get; set; } // Название
    public decimal NewProductPrice { get; set; } // Цена
    public int NewProductStock { get; set; } // Количество
    public string SelectedCategory { get; set; } // Категория

    // Данные покупателя для заказа
    public string CustomerName { get; set; }

    // Выбранные объекты UI
    public Product SelectedProduct { get; set; }
    public CartItem SelectedCartItem { get; set; }
    public Order SelectedOrder { get; set; }

    // Предустановленные категории для выбора
    public ObservableCollection<string> Categories { get; } = new()
    {
        "Ноутбуки", "Комплектующие", "Аксессуары", "Мониторы", "Сеть"
    };

    // Команды для привязки к элементам UI
    public ICommand AddProductCommand { get; }
    public ICommand UpdateProductCommand { get; }
    public ICommand DeleteProductCommand { get; }
    public ICommand AddToCartCommand { get; }
    public ICommand RemoveFromCartCommand { get; }
    public ICommand AddOrderCommand { get; }
    public ICommand DeleteOrderCommand { get; }
    public ICommand ClearFieldsCommand { get; }

    // Конструктор инициализирует команды и загружает данные
    public MainViewModel()
    {
        AddProductCommand = new RelayCommand(async _ => await AddProductAsync(), _ => CanAddProduct());
        UpdateProductCommand = new RelayCommand(async _ => await UpdateProductAsync(), _ => SelectedProduct != null);
        DeleteProductCommand = new RelayCommand(async _ => await DeleteProductAsync(), _ => SelectedProduct != null);
        AddToCartCommand = new RelayCommand(_ => AddToCart(), _ => SelectedProduct != null && SelectedProduct.Stock > 0);
        RemoveFromCartCommand = new RelayCommand(_ => RemoveFromCart(), _ => SelectedCartItem != null);
        AddOrderCommand = new RelayCommand(async _ => await AddOrderAsync(), _ => CartItems.Any() && !string.IsNullOrWhiteSpace(CustomerName));
        DeleteOrderCommand = new RelayCommand(async _ => await DeleteOrderAsync(), _ => SelectedOrder != null);
        ClearFieldsCommand = new RelayCommand(_ => ClearProductFields());

        _ = LoadProductsAsync(); // Загрузка списка товаров
        _ = LoadOrdersAsync();   // Загрузка списка заказов
    }

    private async Task LoadProductsAsync() // Получение всех товаров с сервера
    {
        Products.Clear();
        var products = await _api.GetProductsAsync();
        foreach (var p in products)
            Products.Add(p);
    }

    private async Task LoadOrdersAsync() // Получение всех заказов с сервера
    {
        Orders.Clear();
        var orders = await _api.GetOrdersAsync();
        foreach (var order in orders)
            Orders.Add(order);
    }

    private async Task AddProductAsync() // Добавление нового товара
    {
        var product = new Product
        {
            Name = NewProductName,
            Price = NewProductPrice,
            Stock = NewProductStock,
            Category = SelectedCategory
        };

        await _api.AddProductAsync(product);
        await LoadProductsAsync();
        ClearProductFields();
    }

    private async Task UpdateProductAsync() // Обновление существующего товара
    {
        if (SelectedProduct == null) return;

        SelectedProduct.Name = NewProductName;
        SelectedProduct.Price = NewProductPrice;
        SelectedProduct.Stock = NewProductStock;
        SelectedProduct.Category = SelectedCategory;

        await _api.UpdateProductAsync(SelectedProduct);
        await LoadProductsAsync();
        ClearProductFields();
    }

    private async Task DeleteProductAsync() // Удаление выбранного товара
    {
        if (SelectedProduct == null) return;
        await _api.DeleteProductAsync(SelectedProduct.Id);
        await LoadProductsAsync();
        ClearProductFields();
    }

    private void AddToCart() // Добавление товара в корзину
    {
        var existing = CartItems.FirstOrDefault(x => x.Product.Id == SelectedProduct.Id);
        if (existing != null)
            existing.Quantity++;
        else
            CartItems.Add(new CartItem { Product = SelectedProduct, Quantity = 1 });

        SelectedProduct.Stock--; // Уменьшаем остаток на складе
    }

    private void RemoveFromCart() // Удаление товара из корзины
    {
        if (SelectedCartItem == null) return;
        SelectedCartItem.Product.Stock += SelectedCartItem.Quantity; // Возвращаем остаток
        CartItems.Remove(SelectedCartItem);
        SelectedCartItem = null;
    }

    private async Task AddOrderAsync() // Оформление нового заказа
    {
        var order = new Order
        {
            CustomerName = CustomerName,
            OrderDate = DateTime.Now,
            Items = CartItems.Select(c => new OrderItem
            {
                ProductId = c.Product.Id,
                Quantity = c.Quantity
            }).ToList()
        };

        await _api.AddOrderAsync(order);
        await LoadOrdersAsync();

        CartItems.Clear();
        CustomerName = "";
        await LoadProductsAsync(); // Обновление остатков
    }

    private async Task DeleteOrderAsync() // Удаление существующего заказа
    {
        if (SelectedOrder == null) return;
        await _api.DeleteOrderAsync(SelectedOrder.Id);
        await LoadOrdersAsync();
    }

    private void ClearProductFields() // Очистка полей формы ввода товара
    {
        NewProductName = "";
        NewProductPrice = 0;
        NewProductStock = 0;
        SelectedCategory = null;
    }

    private bool CanAddProduct() => // Проверка валидности данных для добавления товара
        !string.IsNullOrWhiteSpace(NewProductName) &&
        NewProductPrice > 0 &&
        NewProductStock > 0 &&
        !string.IsNullOrWhiteSpace(SelectedCategory);

    public event PropertyChangedEventHandler PropertyChanged; // Событие для INotifyPropertyChanged
    protected void OnPropertyChanged([CallerMemberName] string name = null) =>
        PropertyChanged?.Invoke(this, new PropertyChangedEventArgs(name)); // Уведомление об изменении свойства
}
\end{minted}
	\captionof{listing}{\colorGIT{\href{https://github.com/WebMasterIT/Csharp_Labs/blob/main/6lab/StoreManager_6lab/ViewModels/MainViewModel.cs}{\texttt{MainViewModel.cs}}}}
	\label{lst:main_view_model}
\end{code}
\noindent Описание методов (к листингу~\ref{lst:main_view_model}):
\begin{enumerate}
	\item \texttt{LoadProductsAsync()} — загружает список товаров с сервера.
	\item \texttt{LoadOrdersAsync()} — загружает список заказов с сервера.
	\item \texttt{AddProductAsync()} — формирует товар и отправляет на сервер.
	\item \texttt{UpdateProductAsync()} — обновляет выделенный товар на сервере.
	\item \texttt{DeleteProductAsync()} — удаляет выбранный товар.
	\item \texttt{AddToCart()} — добавляет товар в корзину.
	\item \texttt{RemoveFromCart()} — удаляет товар из корзины.
	\item \texttt{AddOrderAsync()} — создает заказ на основе корзины и отправляет на сервер.
	\item \texttt{DeleteOrderAsync()} — удаляет выбранный заказ.
	\item \texttt{ClearProductFields()} — очищает форму ввода товара.
	\item \texttt{CanAddProduct()} — проверяет, можно ли добавить товар.
\end{enumerate}
Класс \texttt{MainViewModel} является связующим звеном между пользовательским интерфейсом и логикой работы с Web API. Все операции реализованы с учётом асинхронности и привязки данных.

\newpage

\section{\texttt{OrderViewModel.cs} — управление заказами~\texorpdfstring{\iicon{}}{}}
Класс \texttt{OrderViewModel} реализует логику взаимодействия с заказами. Он отвечает за загрузку существующих заказов, создание новых и удаление. Вся логика работает через \texttt{ApiService}, используя команды MVVM.
\begin{code}
	\begin{minted}{csharp}
public class OrderViewModel : BaseViewModel
{
    private readonly ApiService _api = new(); // Сервис API для обмена с сервером

    private Order _selectedOrder; // Выбранный заказ в интерфейсе
    private string _customerName; // Имя покупателя

    public ObservableCollection<Order> Orders { get; } = new(); // Коллекция заказов
    public ObservableCollection<Product> AvailableProducts { get; } = new(); // Доступные для заказа товары
    public ObservableCollection<OrderItem> SelectedItems { get; } = new(); // Позиции текущего заказа

    public string CustomerName
    {
        get => _customerName;
        set => SetProperty(ref _customerName, value); // Установка имени покупателя
    }

    public Order SelectedOrder
    {
        get => _selectedOrder;
        set
        {
            if (SetProperty(ref _selectedOrder, value) && value != null)
            {
                CustomerName = value.CustomerName; // Заполняем имя клиента из заказа
                SelectedItems.Clear(); // Очищаем старые позиции
                foreach (var item in value.Items)
                {
                    SelectedItems.Add(new OrderItem // Копируем позиции
                    {
                        ProductId = item.ProductId,
                        Product = item.Product,
                        Quantity = item.Quantity
                    });
                }
            }
        }
    }

    // Команды для кнопок интерфейса
    public ICommand AddOrderCommand { get; }
    public ICommand DeleteOrderCommand { get; }

    // Конструктор: инициализация и загрузка данных
    public OrderViewModel()
    {
        AddOrderCommand = new RelayCommand(async _ => await AddOrder());
        DeleteOrderCommand = new RelayCommand(async _ => await DeleteOrder(), _ => SelectedOrder != null);

        _ = LoadOrders(); // Загрузка заказов
        _ = LoadAvailableProducts(); // Загрузка товаров
    }

    private async Task LoadOrders() // Получение списка заказов
    {
        Orders.Clear();
        var orders = await _api.GetOrdersAsync();
        foreach (var order in orders)
            Orders.Add(order);
    }

    private async Task LoadAvailableProducts() // Получение списка всех товаров
    {
        AvailableProducts.Clear();
        var products = await _api.GetProductsAsync();
        foreach (var p in products)
            AvailableProducts.Add(p);
    }

    private async Task AddOrder() // Создание нового заказа
    {
        if (string.IsNullOrWhiteSpace(CustomerName) || SelectedItems.Count == 0) return;

        var newOrder = new Order
        {
            CustomerName = CustomerName,
            OrderDate = DateTime.Now,
            Items = SelectedItems.Select(i => new OrderItem
            {
                ProductId = i.Product.Id,
                Quantity = i.Quantity
            }).ToList()
        };

        await _api.AddOrderAsync(newOrder); // Отправка заказа на сервер
        await LoadOrders(); // Обновление списка

        CustomerName = ""; // Очистка формы
        SelectedItems.Clear();
    }

    private async Task DeleteOrder() // Удаление выбранного заказа
    {
        if (SelectedOrder == null) return;

        await _api.DeleteOrderAsync(SelectedOrder.Id);
        Orders.Remove(SelectedOrder);
        SelectedOrder = null;
        SelectedItems.Clear();
        CustomerName = string.Empty;
    }
}
\end{minted}
	\captionof{listing}{\colorGIT{\href{https://github.com/WebMasterIT/Csharp_Labs/blob/main/6lab/StoreManager_6lab/ViewModels/OrderViewModel.cs}{\texttt{OrderViewModel.cs}}}}
	\label{lst:order_view_model}
\end{code}
\noindent Описание методов (к листингу~\ref{lst:order_view_model}):
\begin{enumerate}
	\item \texttt{LoadOrders()} — загружает список всех заказов с сервера.
	\item \texttt{LoadAvailableProducts()} — загружает товары для формирования нового заказа.
	\item \texttt{AddOrder()} — создаёт новый заказ и отправляет его на сервер.
	\item \texttt{DeleteOrder()} — удаляет выбранный заказ.
\end{enumerate}
Класс \texttt{OrderViewModel} выполняет функции управления заказами в клиентской части и полностью опирается на Web API. Реализация соответствует шаблону MVVM и позволяет упростить обработку заказов в интерфейсе WPF.

\newpage

\section{Контрольные вопросы~\texorpdfstring{\iicon{}}{}}
\begin{enumerate}
	\item В чём заключается основное отличие архитектуры лабораторной работы №6 от №4?
	\item Какие преимущества даёт использование ASP.NET Core Web API по сравнению с прямым доступом к базе данных?
	\item Как реализуется добавление товара в лабораторной №6? Чем оно отличается от предыдущей реализации?
	\item Для чего используется класс \texttt{ApiService}? Какие методы он предоставляет?
	\item Как происходит удаление заказа в WPF-приложении, работающем с Web API?
	\item Почему необходимо использовать \texttt{await} при вызовах методов \texttt{ApiService}?
	\item В каком формате передаются данные между клиентом и сервером?
	\item Какие классы относятся к слою ViewModel, и как они связаны с моделью?
	\item Как осуществляется привязка команд в XAML и зачем используются команды MVVM?
	\item Что делает метод \texttt{LoadOrders()} в \texttt{OrderViewModel} и какие данные он загружает?
	\item Какие изменения произошли в механизме сохранения заказов в новой архитектуре?
	\item Почему важно разделение клиентской и серверной логики при разработке приложений?
\end{enumerate}

\end{document}
